\documentclass[12pt]{article}

\usepackage{float}
\usepackage{xcolor}
\usepackage{lmodern}
\usepackage{caption}
\usepackage{booktabs}
\usepackage{listings}
\usepackage{setspace}
\usepackage{graphicx}
\usepackage{geometry}
\usepackage{titlesec}
\usepackage{microtype}
\usepackage[utf8]{inputenc}
\usepackage{amsmath, amssymb, amsfonts}

\geometry{paper=a4paper, layout=a4paper, margin=0.5in}
\setstretch{1.15}

\lstset{
    language=Python,
    basicstyle=\ttfamily\small,
    breaklines=true,
    breakatwhitespace=false,
    columns=fullflexible,
    keepspaces=true,
    showstringspaces=false
}

\begin{document}

\begin{titlepage}
    \begin{center}
        \vspace*{2cm}
        {\Huge \textbf{ ME 4503: Mechanics of Machine}} \\
        \vspace{1cm}
        \Large{\textbf{Assignment II}} \\

        \vfill

        \textit{\large Submitted by:} \\
        \vspace{0.3cm}
        \textbf{Abdullah al Azmi} \\
        \textbf{ID: 220011230} \\

        \vfill

        \textit{\large Submitted to:} \\
        \vspace{0.3cm}
        \textbf{\Large Fatima Tasneem Mohsin (Lecturer)} \\
        \textbf{Department of Mechanical and Production Engineering} \\

        \vfill

        \textbf{Date of Submission: \today}
        \vspace*{1cm}
    \end{center}
\end{titlepage}

\newpage

\section{Question 1: Cam Profile Design}

Here, we have to design a \textit{\textbf{CAM profile}} for \textit\textbf{{Knife Edge follower}} with the CAM specification of \textit{\textbf{50 mm minimum radius, 25 mm shaft diameter, 230 rpm clockwise rotation}}, and the follower specification of \textit{\textbf{43 mm stroke length, SHM (Simple Harmonic Motion) during raise \& fall, 43 mm rise during $120^\circ$ rotation, $30^\circ$ dwell, 43 mm lower during next $90^\circ$, and $120^\circ$ dwell at end}}.

\subsection{Task}

\begin{itemize}

    \item [(a)] Below, a screenshot clearly showing the \textit{\textbf{minimum diameter of the CAM}} is given:

          \begin{figure} [H]
              \centering
              \includegraphics[width=0.6\linewidth,keepaspectratio]{Images/minimum_dia_of_the_shaft}
          \end{figure}

    \item [(b)] Below, a screenshot of the \textit{\textbf{displacement diagram}} is given:

          \begin{figure} [H]
              \centering
              \includegraphics[width=0.6\linewidth,keepaspectratio]{Images/displacement_profile}
          \end{figure}

    \item [(c)] Below, a screenshot of the \textit{\textbf{complete CAM follower system}} is given:

          \begin{figure} [H]
              \centering
              \includegraphics[width=0.6\linewidth,keepaspectratio]{Images/Untitled}
          \end{figure}

    \item [(d)] Below, a screenshot of the \textit{\textbf{Linear Displacement vs. Time Diagram}} is given:

          \begin{figure} [H]
              \centering
              \includegraphics[width=0.6\linewidth,keepaspectratio]{Images/linear_displacement}
          \end{figure}

    \item [(e)] Below, a screenshot of the \textit{\textbf{Linear Velocity vs. Time Diagram}} is given:

          \begin{figure} [H]
              \centering
              \includegraphics[width=0.6\linewidth,keepaspectratio]{Images/linear_velocity}
          \end{figure}

    \item [(f)] Below, a screenshot of the \textit{\textbf{Linear Acceleration vs. Time Diagram}} is given:

          \begin{figure} [H]
              \centering
              \includegraphics[width=0.6\linewidth,keepaspectratio]{Images/linear_acceleration}

          \end{figure}

\end{itemize}

\section{Question 2: Compound Gear Train Design}

Here, we will design a compound gear train that reduces the rotational speed between an input shaft and output shaft. The \textit{\textbf{input shaft}} rotates at \textit{\textbf{1830 rpm}} in the clockwise direction. The \textit{\textbf{output shaft}} rotates at \textit{\textbf{180 rpm}} in the counterclockwise direction. We are considering the following constraints:

\begin{itemize}

    \item [(a)] \textit{No gear in the transmission may act as an idler gear.}
    \item [(b)] \textit{The gear ratio at any individual mesh must not exceed 3:1.}
    \item [(c)] \textit{Available gears have integer tooth numbers ranging from 13 to 85.}
    \item [(d)] \textit{Only one gear exists for each tooth number, meaning the same tooth number cannot be used more than
              once in the system.}
    \item [(e)] \textit{If the difference between the required ratio and calculated ratio is smaller than 0.01, then the gear
              combination can be acceptable.}

\end{itemize}

\subsection{Task}

\begin{itemize}
    \item [(a)] Below, a \textit{\textbf{labeled sketch}} of the gear train configuration is provided:

          \begin{figure}[H]
              \centering
              \includegraphics[width=0.6\linewidth,keepaspectratio]{Images/Figure_1}
          \end{figure}

    \item [(b)] The \textit{\textbf{python code}} used to search for possible solutions:

          \begin{lstlisting}

import pandas as pd

input_rpm = 1830
output_rpm = 180
target_ratio = input_rpm / output_rpm

min_teeth = 13
max_teeth = 85
max_mesh_ratio = 3.0
tolerance = 0.01

def solve_gear_train():
    solutions = []

    for t1 in range(min_teeth, max_teeth + 1):
        for t2 in range(min_teeth, max_teeth + 1):
            r1 = t2 / t1
            if r1 > max_mesh_ratio or t1 == t2: continue

            for t3 in range(min_teeth, max_teeth + 1):
                for t4 in range(min_teeth, max_teeth + 1):
                    for t5 in range(min_teeth, max_teeth + 1):
                        for t6 in range(min_teeth, max_teeth + 1):
                            used = {t1, t2, t3, t4, t5, t6}
                            if len(used) < 6: continue

                            r2 = t4 / t3
                            r3 = t6 / t5
                            if r2 > max_mesh_ratio or r3 > max_mesh_ratio: continue

                            calc_ratio = r1 * r2 * r3
                            if abs(calc_ratio - target_ratio) < tolerance:
                                solutions.append([t1, t2, t3, t4, t5, t6, round(calc_ratio, 4)])
                                if len(solutions) >= 10:
                                    return solutions
    return solutions

columns = ['T1 (In)', 'T2', 'T3', 'T4', 'T5', 'T6 (Out)', 'Calculated Ratio']
results = solve_gear_train()
df = pd.DataFrame(results, columns=columns)

print("\nTop 10 Candidate Gear Combinations:")
print(df.to_string(index=False))
        \end{lstlisting}

    \item [(c)] A \textit{\textbf{table}} of ten candidate gear combinations produced by the program:

          \begin{table}[H]
              \centering
              \begin{tabular}{ccccccc}
                  \textit{\textbf{T1} (In)} & \textit{\textbf{T2}} & \textit{\textbf{T3}} & \textit{\textbf{T4}} & \textit{\textbf{T5}} & \textit{\textbf{T6} (Out)} & \textit{\textbf{Calculated Ratio}} \\
                  13                        & 15                   & 14                   & 42                   & 16                   & 47                         & 10.1683                            \\
                  13                        & 15                   & 16                   & 47                   & 14                   & 42                         & 10.1683                            \\
                  13                        & 15                   & 16                   & 47                   & 17                   & 51                         & 10.1683                            \\
                  13                        & 15                   & 16                   & 47                   & 18                   & 54                         & 10.1683                            \\
                  13                        & 15                   & 16                   & 47                   & 19                   & 57                         & 10.1683                            \\
                  13                        & 15                   & 16                   & 47                   & 20                   & 60                         & 10.1683                            \\
                  13                        & 15                   & 16                   & 47                   & 21                   & 63                         & 10.1683                            \\
                  13                        & 15                   & 16                   & 47                   & 22                   & 66                         & 10.1683                            \\
                  13                        & 15                   & 16                   & 47                   & 23                   & 69                         & 10.1683                            \\
                  13                        & 15                   & 16                   & 47                   & 24                   & 72                         & 10.1683                            \\
              \end{tabular}
          \end{table}

    \item [(d)] The selected optimal gear set is the first one: \textit{\textbf{T1(In) = 13, T2 = 15, T3 = 14, T4 = 42, T5 = 16 \& T6 (Out) = 47}}. The reasons for being chosen the first combination is:

          \begin{itemize}

              \item [(1)] A compound gear train with three stages (three meshes) reverses the direction of rotation an odd number of times.

                    \begin{itemize}
                        \item [(a)] \textit{Stage 1: CW $\rightarrow$ CCW Stage}
                        \item [(b)] \textit{Stage 2: CCW $\rightarrow$ CW Stage}
                        \item [(c)] \textit{Stage 3: CW $\rightarrow$ CCW Stage}
                    \end{itemize}

                    Since the requirement is for the output shaft to rotate counterclockwise, a 3-stage design is necessary and correct.

              \item [(2)] The first set is superior to others for two reasons:

                    \begin{itemize}

                        \item [(a)] The maximum gear size used is \textit{\textbf{47 teeth}} where the later design used gears up to \textit{\textbf{72 teeth}}. Smaller gears result in a more compact gearbox, lower weight, and lower material cost.

                        \item [(b)] Using gear counts like \textit{\textbf{13, 47, and 17}} is generally preferred in mechanical design. These are often prime or have fewer common factors, which ensures that the same teeth don't meet as frequently.

                    \end{itemize}

              \item [(3)] The ratios for this set are:

                    \begin{itemize}
                        \item [(a)] $R_1$ = 15/13 $\approx$ 1.15
                        \item [(b)] $R_2$ = 42/14 $\approx$ 3.00
                        \item [(c)] $R_3$ = 47/16 $\approx$ 2.94
                    \end{itemize}

                    All individual meshes are $\leq 3.0$, satisfying constraint \textit{\textbf{(b)}}. While $R_1$ is quite small, keeping the largest gear to 47 teeth (instead of 72) is a significant advantage for the overall footprint of the assembly.

              \item [(4)] All gears are part of a mesh that contributes to the ratio. So, there are \textit{no idlers}.

              \item [(5)]  All tooth numbers $\{13, 15, 14, 42, 16, 47\}$ are unique, satisfying constraint \textit{\textbf{(d)}}.

              \item [(6)] All values are within the $[13, 85]$ range.

          \end{itemize}
\end{itemize}

\end{document}
